\documentclass[]{article}

\usepackage{amsmath}
\usepackage{multirow}
\usepackage{natbib}
\usepackage{graphicx} 


\begin{document}

Cosmological N-body simulations are computationally demanding, which limits the generation of large statistical ensembles required for precision cosmology. We address this challenge by developing and benchmarking a GPU-accelerated Particle-Mesh (PM) simulation using the NVIDIA Warp framework. Our implementation leverages custom Warp kernels for Cloud-in-Cell (CIC) particle-grid operations and the cuFFT library for the Poisson solver. When simulating $512^3$ particles, our GPU code achieves a 1,348-fold per-step speedup over an equivalent multi-threaded CPU implementation, with the largest gains (over 2,000$\times$) in the particle-grid interaction steps. To validate the physical fidelity of the pipeline's components, we generated initial conditions at z=127 using the Zel'dovich Approximation. The resulting matter power spectrum, when linearly evolved to z=0, agrees with theoretical predictions from CAMB to within 4\% on large scales (k < 0.04 h/Mpc), while deviations on smaller scales are consistent with the known limitations of the approximation used. This work demonstrates that NVIDIA Warp is a highly effective tool for building performant cosmological codes, enabling the rapid generation of mock universes for large-scale structure analysis.

\end{document}
                